\documentclass[11pt,letterpaper]{article}

% ── Geometry ──────────────────────────────────────────────────────────────────
\usepackage[margin=1in]{geometry}

% ── Pagination ────────────────────────────────────────────────────────────────
\usepackage[all]{nowidow}    % Require ≥2 lines on each side of a page break
\brokenpenalty=10000         % No hyphenated words across page breaks
\raggedbottom                % Allow uneven page bottoms rather than stretching

% ── Fonts ─────────────────────────────────────────────────────────────────────
\usepackage[T1]{fontenc}
\usepackage{newpxtext}     % Palatino-based serif (TeX Gyre Pagella)
\usepackage{newpxmath}     % Matching math font

% ── Colors ────────────────────────────────────────────────────────────────────
\usepackage[dvipsnames]{xcolor}
\definecolor{SectionBlue}{HTML}{1B3A5C}
\definecolor{AccentGray}{HTML}{4A4A4A}
\definecolor{QuoteBorder}{HTML}{8B9DAF}
\definecolor{RiskRed}{HTML}{C0392B}
\definecolor{RiskAmber}{HTML}{E67E22}
\definecolor{RiskGreen}{HTML}{27AE60}
\definecolor{BoxBg}{HTML}{F7F8FA}

% ── Tables ────────────────────────────────────────────────────────────────────
\usepackage{booktabs}
\usepackage{tabularx}

% ── Boxes ─────────────────────────────────────────────────────────────────────
\usepackage[framemethod=tikz]{mdframed}

% ── Lists ─────────────────────────────────────────────────────────────────────
\usepackage{enumitem}
\usepackage{needspace}

% ── Hyperlinks ────────────────────────────────────────────────────────────────
\usepackage{hyperref}
\hypersetup{
  colorlinks=true,
  linkcolor=SectionBlue,
  urlcolor=SectionBlue,
  citecolor=SectionBlue,
  pdfauthor={Punt Labs},
  pdftitle={PR/FAQ: punt-kit},
  pdfsubject={Working Backwards — Product Discovery},
}

% ── Bibliography ─────────────────────────────────────────────────────────────
\usepackage[style=numeric,backend=biber,sorting=none]{biblatex}
\addbibresource{prfaq.bib}

% ── Section styling ───────────────────────────────────────────────────────────
\usepackage{titlesec}
\titleformat{\section}
  {\needspace{6\baselineskip}\Large\bfseries\color{SectionBlue}}
  {}
  {0pt}
  {}

\titleformat{\subsection}
  {\needspace{5\baselineskip}\large\bfseries\color{AccentGray}}
  {}
  {0pt}
  {}

\titleformat{\subsubsection}
  {\normalsize\bfseries\color{AccentGray}}
  {}
  {0pt}
  {}
\titlespacing*{\subsubsection}{0pt}{1em}{0.5em}[0pt]

% ── Document stage ───────────────────────────────────────────────────────────
\newcommand{\prfaqstage}[1]{\def\theprfaqstage{#1}}
\prfaqstage{internal-adoption}

% ── Document version ─────────────────────────────────────────────────────────
\newcommand{\prfaqversion}[2]{\def\theprfaqmajor{#1}\def\theprfaqminor{#2}}
\prfaqversion{2}{0}

% ── Header/footer ────────────────────────────────────────────────────────────
\usepackage{fancyhdr}
\pagestyle{fancy}
\fancyhf{}
\fancyhead[L]{\footnotesize\textcolor{AccentGray}{Stage: \theprfaqstage{} \enspace\textbar\enspace v\theprfaqmajor.\theprfaqminor}}
\fancyhead[R]{\footnotesize\textcolor{AccentGray}{\thepage}}
\renewcommand{\headrulewidth}{0pt}

% ══════════════════════════════════════════════════════════════════════════════
% Custom environments
% ══════════════════════════════════════════════════════════════════════════════

% ── \prsection{title} — styled press release section header ──────────────────
\newcommand{\prsection}[1]{%
  \vspace{1em}%
  {\large\bfseries\color{SectionBlue}#1}%
  \vspace{0.3em}\par%
}

% ── customerquote — indented, italic customer quote ──────────────────────────
\newmdenv[
  topline=false,
  bottomline=false,
  rightline=false,
  linewidth=3pt,
  linecolor=QuoteBorder,
  backgroundcolor=BoxBg,
  innerleftmargin=12pt,
  innerrightmargin=12pt,
  innertopmargin=8pt,
  innerbottommargin=8pt,
  skipabove=\baselineskip,
  skipbelow=\baselineskip,
]{customerquote}

% ── spokespersonquote — indented spokesperson quote ──────────────────────────
\newmdenv[
  topline=false,
  bottomline=false,
  rightline=false,
  linewidth=3pt,
  linecolor=SectionBlue,
  backgroundcolor=BoxBg,
  innerleftmargin=12pt,
  innerrightmargin=12pt,
  innertopmargin=8pt,
  innerbottommargin=8pt,
  skipabove=\baselineskip,
  skipbelow=\baselineskip,
]{spokespersonquote}

% ── faqpair — numbered Q&A pair with bold question ─────────────────────────
\newcounter{faqnum}
\newenvironment{faqpair}[1]{%
  \needspace{4\baselineskip}%
  \vspace{0.5em}%
  \refstepcounter{faqnum}%
  \noindent\textbf{\color{SectionBlue}Q\thefaqnum: #1}\par\nopagebreak[4]%
  \vspace{0.2em}%
  \begin{adjustwidth}{1em}{0pt}%
  \setlength{\parindent}{0pt}%
}{%
  \end{adjustwidth}%
  \vspace{0.5em}%
}
% ── \faqref{faq:slug} — clickable "FAQ 3" reference ──────────────────────────
\newcommand{\faqref}[1]{\hyperref[#1]{FAQ~\ref*{#1}}}

% ── \featureitem — numbered, referenceable feature entry ──────────────────────
\newcounter{featurenum}
\newcommand{\featureitem}[2]{%
  \refstepcounter{featurenum}%
  \item[\textbf{\color{SectionBlue}F\thefeaturenum.}]%
  \textbf{#1} --- #2%
}
\newcommand{\featureref}[1]{\hyperref[#1]{Feature~\ref*{#1}}}

% ── adjustwidth (for faqpair indentation) ────────────────────────────────────
\usepackage{changepage}

% ══════════════════════════════════════════════════════════════════════════════
% Document
% ══════════════════════════════════════════════════════════════════════════════

\begin{document}

% ── Headline block ───────────────────────────────────────────────────────────
\begin{center}
  {\LARGE\bfseries\color{SectionBlue}
  punt-kit Is the Kit Manager for the\\Punt Labs Org} \\[0.8em]
  {\large\color{AccentGray}
  One tool manages the org's standards, its tools, its workflows, and its
  compliance checks --- for the operator and for the AI agents that do the
  work across roughly thirty repositories} \\[0.3em]
  {\large\color{AccentGray}
  \texttt{/punt:auto} turns recurring prompts into reviewed programs:
  playbooks that mix deterministic shell with LLM judgment, and --- next ---
  prompt state machines that drive the agent instead of the operator
  retyping the loop}
\end{center}

\vspace{0.5em}

% ══════════════════════════════════════════════════════════════════════════════
% PRESS RELEASE
% ══════════════════════════════════════════════════════════════════════════════

\section*{Press Release}

PUNT LABS --- December 1, 2026 --- punt-kit is the kit manager for the Punt
Labs org. One tool now manages the four things the org runs on: its
\textbf{standards} (author them, activate them in a repo, audit against them),
its \textbf{tools} (install them, enable them, wire their guidance into a
session), its \textbf{workflows} (the release engine and the playbooks), and
its \textbf{compliance} checks (audit, PII scan). The customer is Punt Labs
itself --- one operator and the AI agents that do the work across roughly
thirty repositories. External adoption is not the goal of this version and is
explicitly deferred (see \faqref{faq:what-is-it}, \faqref{faq:external}).

The new capability in this release is \texttt{/punt:auto} as a
\textbf{self-prompting engine}: the recurring prompts the operator used to type
by hand become playbooks --- reviewed, versioned programs that generate the
prompts and drive the agent themselves (see \faqref{faq:self-prompting}).

\prsection{Problem}

Punt Labs is one operator and a roster of AI agents working roughly thirty
repositories. Two costs dominate, and both scale with the number of repos and
sessions rather than with the amount of code.

The first is context. Every agent session starts cold, with no memory of what
came before~\cite{anthropic2025contextengineering}, and the org's conventions
--- how a release runs, what a README must contain, when a PR may merge ---
live in prose that each session has to find. A study of 253 CLAUDE.md files
found these manifests are optimized for code execution rather than
organizational convention: only 8.7\% include security
guidance~\cite{arxiv2509.14744}. Static context files also have no feedback
loop, so nothing notices when a repo has drifted from the document that
describes it~\cite{claudecodememory}. Catching that drift in review is the
expensive path: rote standards checking is estimated at roughly four hours per
engineer per week~\cite{deepsource2025staticanalysis}.

The second is the operator typing the same prompt again. Watching a pull
request to merge, sweeping the backlog, running a release --- these are loops
with named states and guarded transitions, and today the operator re-authors
each one as a fresh cron prompt, once per PR. A hand-typed loop is
unreviewable, unversioned, and evaporates when the session compacts. It is the
same toil the standards corpus removed for conventions, still untouched for
procedure (see \faqref{faq:self-prompting}).

\prsection{Solution}

\textbf{Standards.} The corpus is 29 normative prose documents in
\texttt{standards/}, plus per-language coding rules in \texttt{lang-rules/}
that agents auto-load when they touch a matching file. They are prose on
purpose: the consumer is a language model and a human reviewer, and neither
needs the rules ground into assertions first. What is machine-checkable is
checked separately and deterministically --- \texttt{punt audit} is
hand-written Python that verifies the org's conventions in a repo (CI
workflows, tool config, permissions, Makefile targets, README SHA pins, plugin
version sync, branch-protection rulesets) and reports each finding with its
rule. \texttt{punt init} scaffolds a new repo that already satisfies them.

\textbf{Tools.} Punt Labs ships around a dozen of its own tools, and a repo
should be able to say which ones it uses. The
\texttt{standards/tool-enable-disable.md} convention is how: each tool owns a
\texttt{.punt-labs/<tool>/} subtree, an \texttt{enabled} marker records the
choice, and the host \texttt{CLAUDE.md} carries exactly one
\texttt{@}-import line per enabled tool so its guidance composes into every
session. \texttt{enable} and \texttt{disable} are the verbs; \texttt{disable}
is non-destructive. This is the kit-manager direction --- punt-kit installs,
enables, and audits the kit (see \faqref{faq:kit-manager}).

\textbf{Workflows.} \texttt{punt release} is an 11-phase release engine ---
preflight, version bump, build, release PR, tag, CI wait, GitHub release, PyPI
verification, post-release, cross-repo propagation, and a final verify pass.
It is the most-used feature in the org and runs on every release.

\textbf{Compliance.} \texttt{punt audit} and \texttt{punt pii} are the
deterministic checks: convention conformance and a scan for emails, home
directory paths, and hostnames, with a \texttt{--staged} mode for pre-commit.

\textbf{The engine underneath.} \texttt{/punt:auto} runs playbooks: YAML
programs whose steps are either deterministic shell or scoped LLM judgment,
with postconditions and a per-step failure strategy. The executor is an LLM
following \texttt{skills/auto/SKILL.md} --- not a daemon --- which is what lets
a judgment step sit next to a shell step in one program. The release engine is
the shipped proof: eleven deterministic phases with LLM diagnosis on failure.
The drafted next step turns the DSL into a prompt state machine --- named
states, guarded transitions, poll ticks, and persisted run-state with
crash/resume --- so a loop is a reviewed file rather than a retyped prompt
(see \faqref{faq:self-prompting}).

\prsection{Customer Quote}

\begin{customerquote}
\textit{``I run the releases, and every one of them goes through
\texttt{/punt:auto release} --- eleven phases, and when a phase fails the
executor reads the error and tries a fix instead of handing me a stack trace.
That is the part I no longer think about. The part I still think about is the
watching: I hand-type the same PR poll prompt every time a PR opens, and I
typed it twice in one session last week. Those prompts are a state machine I
re-derive from memory each time. I want to run the file instead.''}

\raggedleft --- \textbf{Claude Agento}, COO/VP Engineering agent, Punt Labs
\end{customerquote}

\prsection{Getting Started}

Install with \texttt{install.sh} (or \texttt{uv tool install punt-kit} if uv is
already present). In a new repo, \texttt{punt init} scaffolds CI workflows,
tool config, permissions, and a \texttt{CLAUDE.md} that references the
standards. \texttt{punt audit} reports where an existing repo diverges.
\texttt{/punt:auto list} shows the playbooks available, and
\texttt{/punt:auto release} runs a release end to end. Nothing phones home;
every command runs locally against the checkout in front of it.

\prsection{Spokesperson Quote}

\begin{spokespersonquote}
``The first version of this document described a tool that read our markdown
standards and extracted checkable assertions from them. We never built that,
and after two years of using the thing, I do not want it. The consumer of a
standard is a model and a reviewer, and prose is the right format for both;
the checks that can be deterministic are better written as code than derived
from a paragraph. What punt-kit actually became is the manager for the org's
kit --- the standards, the tools, the workflows, the compliance checks --- and
the honest version of this document says so. The part I am investing in is
\texttt{/punt:auto}. Every recurring prompt I type by hand is a program I have
not written yet. Turning those into playbooks, and then into state machines
with guards and resume, means the tool writes the prompt and the agent runs
the loop. That is the leverage. Whether anyone outside Punt Labs wants this is
a question for later, deliberately.''

\raggedleft --- \textbf{J. Freeman}, Creator, Punt Labs
\end{spokespersonquote}

\prsection{Call to Action}

punt-kit is public at {\small\texttt{github.com/punt-labs/punt-kit}}, on PyPI
as \texttt{punt-kit}, and on the punt-labs Claude Code plugin marketplace as
\texttt{punt@punt-labs}, under the MIT license. It is public because the org's
other tools install it, not because it is being marketed. For a Punt Labs repo
that has not adopted it: run \texttt{punt audit}, fix what it reports, and put
the standards reference in \texttt{CLAUDE.md}.

\newpage

% ══════════════════════════════════════════════════════════════════════════════
% FREQUENTLY ASKED QUESTIONS
% ══════════════════════════════════════════════════════════════════════════════

\section*{Frequently Asked Questions}

% ── External FAQs (Customer-facing) ──────────────────────────────────────────

\subsection*{External FAQs}

\begin{faqpair}{What is punt-kit and who is it for?}\label{faq:what-is-it}
  punt-kit is a development lifecycle engine that reads engineering standards
  from a directory of markdown documents and enforces them as executable checks.
  It is designed for engineering leads and platform-minded engineers at teams of
  5--50 people who maintain multiple repositories and want their standards
  enforced automatically --- not just documented. The engine projects through
  four surfaces: a CLI (\texttt{punt}) for terminal use, a Claude Code plugin
  (\texttt{/punt}) for AI coding sessions, an MCP server so any AI agent can
  call the same checks, and a REST API for CI pipelines and dashboards
  (see \faqref{faq:projection}).
\end{faqpair}

\begin{faqpair}{How is this different from projen, Backstage, or Trunk?}\label{faq:how-different}
  Each of these tools solves a narrower problem:
  \begin{itemize}[nosep]
    \item \textbf{projen}~\cite{projen2024github} enforces project
      configuration through TypeScript project types. It is the strongest prior
      art for ``standards as code across repos.'' punt-kit differs in three
      ways: it reads standards from existing markdown rather than requiring
      TypeScript; it integrates with AI coding agent sessions as a first-class
      citizen; and it targets teams that already have standards in docs, not
      greenfield CDK infrastructure teams.
    \item \textbf{Backstage}~\cite{backstage2025fiveyears} is a developer
      portal for large enterprises (3,000+ adopters, but adoption stalls below
      10\% in most non-Spotify organizations). It requires a dedicated
      full-time team to run. punt-kit targets 5--50 engineer teams that cannot
      staff a platform team.
    \item \textbf{Trunk}~\cite{trunkio2024codequality} is a metalinter that
      orchestrates 100+ linting tools with unified output. It covers the
      linting layer. punt-kit covers the layer above it: project structure, CI
      templates, release pipelines, naming conventions, and cross-repo rollout.
  \end{itemize}
  No existing tool unifies audit, scaffold, rollout, and release under a single
  standards-as-data model with AI agent integration
  (see \faqref{faq:competitors}).
\end{faqpair}

\begin{faqpair}{How do I get started?}
  Install with \texttt{pip install punt-kit} or
  \texttt{uv tool install punt-kit}. Run \texttt{punt audit} in a repository.
  Review the findings. Customize the standards directory if the defaults do not
  match your team's conventions. Most teams see their first compliance report
  within five minutes.
\end{faqpair}

\begin{faqpair}{How much does it cost?}
  punt-kit is free and open-source under the MIT license. There are no paid
  tiers, no usage limits, and no telemetry. The tool runs entirely on your
  machine; it does not phone home (see \faqref{faq:revenue-model}).
\end{faqpair}

\begin{faqpair}{What happens to my data?}
  punt-kit runs locally. It reads files from your repositories and your
  standards directory. It does not transmit data to any external server. The MCP
  server and REST API run on your machine (or your infrastructure); findings
  are never sent elsewhere. When used as a Claude Code plugin, the standards
  context is provided to the AI agent session but never stored by punt-kit
  itself.
\end{faqpair}

\begin{faqpair}{Why does punt-kit have four surfaces (CLI, plugin, MCP, API)?}\label{faq:projection}
  Standards enforcement only works if every actor in the development loop uses
  the same checks. A human typing in a terminal needs a CLI. An AI coding agent
  in Claude Code needs a plugin. A different AI agent (Cursor, Windsurf,
  Codex) needs an MCP server it can call. A CI pipeline or internal dashboard
  needs a REST API. All four call the same engine --- the checks, the standards,
  and the findings are identical regardless of which surface invoked them. This
  projection pattern means adopting punt-kit does not lock a team into any
  single agent or workflow.
\end{faqpair}

\begin{faqpair}{Can I use punt-kit with my own standards instead of the defaults?}
  Yes. The engine and the standards pack are separable. punt-kit ships with the
  punt-labs conventions as the default, but any team can point it to their own
  standards directory --- a local folder, a git repository, or a URL. The engine
  reads the rules; it does not decide what ``good'' looks like
  (see \featureref{feat:custom-standards}).
\end{faqpair}

% ── Internal FAQs (Business-facing) ──────────────────────────────────────────

\newpage
\subsection*{Internal FAQs}

\subsubsection*{Value \& Market}

\begin{faqpair}{What is the total addressable market?}\label{faq:tam}
  \textbf{Bottom-up estimate (with caveats):}

  The platform engineering services market was valued at \$7.19 billion in 2024
  and is projected to reach \$40.17 billion by 2032, growing at 23.99\%
  CAGR~\cite{snsconsulting2024platformmarket}. Gartner predicts that 80\% of
  large software engineering organizations will have platform engineering teams
  by 2026, up from 45\% in 2022~\cite{gartnerplatformeng2022}. Google Cloud
  found that 55\% of global organizations have already adopted platform
  engineering, with 86\% saying it is essential to realizing the full business
  value of AI~\cite{googlecloud2024platformengineering}.

  \textbf{Honest gap:} These numbers describe large organizations with
  dedicated platform teams --- the opposite of punt-kit's target segment.
  No public data source directly measures how many 5--50 engineer teams
  maintain multiple repos with codified standards. GitHub hosts 3+ million
  organizations, but there is no breakdown by team size or standards maturity.

  \textbf{Hypothesis:} Smaller teams have the same standards-enforcement pain
  but cannot afford Backstage or a dedicated platform engineer. punt-kit's
  bet is that this segment is underserved precisely because existing solutions
  require enterprise-scale investment. This hypothesis is unvalidated and is the
  highest-priority item to test (see \faqref{faq:next-step}).
\end{faqpair}

\begin{faqpair}{What evidence do we have that customers want this?}\label{faq:customer-evidence}
  \textbf{At hypothesis stage, we have no primary customer evidence.} What we
  have instead:

  \begin{itemize}[nosep]
    \item \textbf{Dogfooding:} punt-kit was built to solve a real problem
      across 14 repositories in the punt-labs org. The tool exists because the
      manual alternative was unsustainable.
    \item \textbf{Academic evidence:} Four arXiv papers from 2025--2026
      validate the context problem. A study of 283 development sessions found
      that single-file agent manifests ``do not scale beyond modest
      codebases''~\cite{arxiv2602.20478}. A study of 2,303 context files found
      security specified in only 14.5\%~\cite{arxiv2511.12884}. Agent sessions
      without context files consume 28.64\% more runtime and 16.58\% more
      tokens~\cite{arxiv2601.20404}.
    \item \textbf{Industry signals:} Code review consumes an estimated 4 hours
      per week per engineer on standards enforcement that could be
      automated~\cite{deepsource2025staticanalysis}. 88\% of developers cite
      negative impacts from AI-generated code~\cite{sonar2026stateofcode}. 23\%
      of developers use AI agents at least weekly and that number is
      growing~\cite{stackoverflow2025survey}.
    \item \textbf{Competitor existence:} projen, Trunk, Copier, and Backstage
      each solve a slice of this problem, confirming market demand for at least
      partial solutions.
  \end{itemize}

  The missing evidence is direct conversations with 5--50 engineer teams
  outside punt-labs confirming that they prioritize this problem highly enough
  to adopt a new tool.
\end{faqpair}

\begin{faqpair}{Who are the competitors and why will we win?}\label{faq:competitors}
  \renewcommand{\arraystretch}{1.4}
  \begin{tabularx}{\linewidth}{@{} l X X @{}}
  \toprule
  \textbf{Tool} & \textbf{What It Does Well} & \textbf{Gap punt-kit Fills} \\
  \midrule
  projen & Standards-as-code for project config across N
    repos~\cite{projen2024github} & Requires TypeScript; AWS-centric; no AI
    agent integration; high learning curve \\
  Backstage & Developer portal with golden path
    scaffolding~\cite{backstage2025fiveyears} & Requires dedicated team to run;
    stalls at under 10\% adoption outside Spotify \\
  Trunk & Metalinter orchestrating 100+
    tools~\cite{trunkio2024codequality} & Covers linting only; no scaffolding,
    rollout, or release management \\
  Copier & Template scaffolding with lifecycle
    updates~\cite{copier2024comparisons} & No enforcement/audit; no AI
    integration; template-only \\
  MegaLinter & Open-source CI linting orchestrator & CI-focused; no scaffolding
    or release; similar scope to Trunk \\
  \bottomrule
  \end{tabularx}
  \renewcommand{\arraystretch}{1.0}

  \medskip
  punt-kit's structural advantage is the combination of four properties no
  competitor has together: (1) standards from existing markdown, not new
  TypeScript or YAML; (2) unified lifecycle (audit + scaffold + rollout +
  release) under one data model; (3) a projection architecture that serves
  CLI, plugin, MCP, and REST from the same engine; (4) agent-agnostic
  integration --- any tool that speaks MCP can enforce standards, not just
  Claude Code. If a competitor copies one of these, the remaining three still
  differentiate.
\end{faqpair}

\begin{faqpair}{What is the customer acquisition strategy?}
  \textbf{Primary channel:} Claude Code plugin marketplace. punt-kit is already
  published on the punt-labs marketplace. Engineers who use Claude Code discover
  \texttt{/punt audit} in their workflow and adopt the CLI for terminal use.

  \textbf{Secondary channels:}
  \begin{itemize}[nosep]
    \item PyPI discovery (\texttt{pip install punt-kit})
    \item GitHub search and stars (open-source visibility)
    \item Blog posts and conference talks on ``standards as executable
      assertions''
    \item Word of mouth from dogfooding across punt-labs projects
  \end{itemize}

  \textbf{Conversion assumption:} Free, open-source, no signup friction.
  Adoption is limited by awareness and by whether the problem is painful enough
  to motivate a tool switch, not by pricing barriers.
\end{faqpair}

\begin{faqpair}{What is the next step to validate the vision?}\label{faq:next-step}
  Interview 15--20 engineering leads at Series~A/B companies with 5--50
  engineers and 5+ repositories. The interview should answer three questions:
  \begin{enumerate}[nosep]
    \item Is standards enforcement a top-3 pain point for you?
    \item How do you currently propagate standard changes across repos?
    \item Would you adopt a CLI that audits and fixes standards compliance
      automatically?
  \end{enumerate}
  If fewer than 8 of 20 interviewees identify this as a top-3 pain, the target
  segment hypothesis is wrong and we should either broaden the customer
  definition or narrow to a more specific persona.
\end{faqpair}

\subsubsection*{Technical}

\begin{faqpair}{What are the major technical risks?}
  \begin{enumerate}[nosep]
    \item \textbf{Standards parsing ambiguity.} Markdown is not a schema
      language. Extracting checkable assertions from prose requires either
      structured conventions in the markdown or LLM-assisted parsing. Mitigation:
      the engine uses deterministic parsing first and falls back to LLM only for
      genuinely ambiguous cases (see \featureref{feat:deterministic-audit}).
    \item \textbf{Cross-repo rollout safety.} Automatically modifying files
      across N repos risks breaking things at scale. Mitigation: dry-run mode
      previews every change before applying; rollout never force-pushes or
      merges without green CI (see \featureref{feat:rollout}).
    \item \textbf{Claude Code plugin API stability.} The plugin hooks and slash
      commands depend on Claude Code's extension API, which is evolving.
      Mitigation: punt-kit already ships a working plugin (v0.2.0) and tracks
      API changes as part of ongoing development.
  \end{enumerate}
\end{faqpair}

\begin{faqpair}{What dependencies exist on other teams or systems?}
  punt-kit depends on Claude Code for plugin distribution and on PyPI for CLI
  distribution. Both are stable, public infrastructure. There are no
  dependencies on internal services, proprietary APIs, or third-party SaaS.
  The LLM-powered reconciliation feature (\texttt{/punt reconcile}) requires an
  Anthropic API key, but all deterministic features work without one.
\end{faqpair}

\begin{faqpair}{What is the estimated development timeline?}
  punt-kit v0.1.0 shipped February 2026 with \texttt{/punt audit},
  \texttt{/punt reconcile}, \texttt{/punt init}, and \texttt{/punt pii}. The
  vision document describes ten additional commands. Phased delivery:
  \begin{itemize}[nosep]
    \item \textbf{Phase 1 (shipped):} Audit, reconcile, init, PII scan
    \item \textbf{Phase 2 (Q2 2026):} Preflight, per-standard audit
      subcommands, doctor
    \item \textbf{Phase 3 (Q3 2026):} Release automation, changelog, status
    \item \textbf{Phase 4 (Q4 2026):} Cross-repo rollout, PR workflow
    \item \textbf{Phase 5 (2027):} Autopilot (autonomous backlog execution)
  \end{itemize}
  If timeline compresses, Phase 5 is cut. Phase 2 is the minimum viable next
  milestone; it validates the ``standards as data'' engine pattern with
  deterministic, per-rule checks.
\end{faqpair}

\subsubsection*{Business}

\begin{faqpair}{What is the revenue model?}\label{faq:revenue-model}
  punt-kit is free and open-source with no planned paid tiers. The strategic
  value is:
  \begin{itemize}[nosep]
    \item \textbf{Ecosystem presence:} punt-kit establishes Punt Labs as a
      credible voice in developer tooling, driving awareness for adjacent
      products (quarry, biff).
    \item \textbf{Dogfooding infrastructure:} punt-kit enforces standards
      across all punt-labs projects, reducing maintenance cost for the org's own
      tools.
    \item \textbf{Platform engineering credibility:} The standards-as-data
      approach, if validated, positions Punt Labs to offer consulting or
      enterprise features later. But monetization is explicitly not the goal
      for v1.
  \end{itemize}
\end{faqpair}

\begin{faqpair}{What are the key metrics for success?}
  \begin{enumerate}[nosep]
    \item \textbf{Adoption:} 50 GitHub stars and 5 non-punt-labs organizations
      using punt-kit within 6 months of v1.0. Measurement: GitHub stars, PyPI
      download counts, Claude Code plugin install count.
    \item \textbf{Retention:} Of adopting orgs, 3+ run \texttt{punt audit}
      at least weekly after 3 months.
    \item \textbf{Standards pack diversity:} At least 2 non-punt-labs standards
      packs contributed or published within 12 months (evidence that the
      engine/pack separation works).
    \item \textbf{Decision threshold:} If after 6 months, fewer than 3 non-punt-labs
      organizations are using punt-kit actively, revisit the customer
      definition and consider whether the tool should remain internal-only.
  \end{enumerate}
\end{faqpair}

\begin{faqpair}{Why now? What has changed?}\label{faq:why-now}
  Three shifts converged:
  \begin{enumerate}[nosep]
    \item \textbf{AI agents are generating more code faster.} 23\% of
      developers use AI agents at least weekly~\cite{stackoverflow2025survey},
      and developers use AI in 60\% of their
      work~\cite{anthropic2026agentictrendsreport}. More code, generated
      faster, with no awareness of organizational conventions means standards
      enforcement scales worse than ever.
    \item \textbf{Agent context infrastructure is maturing.} CLAUDE.md,
      AGENTS.md, and Claude Code plugins now provide hooks to inject standards
      context into every agent session. This infrastructure did not exist 18
      months ago.
    \item \textbf{Platform engineering is going mainstream but tools assume
      enterprise scale.} 55\% of organizations have adopted platform
      engineering~\cite{googlecloud2024platformengineering}, but the leading
      tools (Backstage, projen) require dedicated teams. Small teams need
      the same enforcement without the overhead.
  \end{enumerate}
\end{faqpair}

% ══════════════════════════════════════════════════════════════════════════════
% FOUR RISKS ASSESSMENT
% ══════════════════════════════════════════════════════════════════════════════

\newpage
\section*{Risk Assessment}

\renewcommand{\arraystretch}{1.6}
\begin{tabularx}{\textwidth}{@{} l l X @{}}
\toprule
\textbf{\color{SectionBlue}Risk} & \textbf{\color{SectionBlue}Rating} & \textbf{\color{SectionBlue}Assessment} \\
\midrule
Value & \textcolor{RiskAmber}{Medium} &
  The problem is real and well-evidenced by academic research and industry
  signals. The risk is whether 5--50 engineer teams consider standards
  enforcement a top-3 pain point worth adopting a new tool for. No primary
  customer evidence exists outside dogfooding. The ``Why now'' argument
  (AI agents amplifying drift) is strong but untested with the target
  segment. \\
Usability & \textcolor{RiskGreen}{Low} &
  Three-step onboarding with zero configuration. The tool fits into existing
  terminal and Claude Code workflows rather than replacing them. Time-to-value
  is under five minutes. The main usability risk is standards parsing:
  if teams' existing markdown documents are too unstructured for the engine
  to parse, the setup cost increases. \\
Feasibility & \textcolor{RiskGreen}{Low} &
  Core technology is proven. v0.1.0 is shipped and running across 14
  repositories. The team has domain expertise (the tool was built to solve
  the team's own problem). The hardest unsolved problem is cross-repo rollout
  safety, which is deferred to Phase 4. No dependencies on unproven
  technology. \\
Viability & \textcolor{RiskAmber}{Medium} &
  punt-kit is free with no direct revenue plan. Strategic value (ecosystem
  presence, dogfooding infrastructure) is real but hard to measure.
  The risk is opportunity cost: time spent on punt-kit is time not spent on
  revenue-generating products. Mitigated by the fact that punt-kit directly
  reduces maintenance cost across all other punt-labs projects. \\
\bottomrule
\end{tabularx}
\renewcommand{\arraystretch}{1.0}

% ══════════════════════════════════════════════════════════════════════════════
% FEATURE APPENDIX
% ══════════════════════════════════════════════════════════════════════════════

\newpage
\section*{Feature Appendix}

Defines the scope boundary for the product. Every feature is explicitly
categorized to prevent scope creep and force prioritization decisions upfront.

\subsection*{Must Do}

Features that are essential for launch. Without these, the product does not
solve the core problem.

\begin{enumerate}[nosep,leftmargin=2.5em]
  \featureitem{Standards-as-data engine}{reads a directory of markdown
    documents, extracts checkable assertions, and runs them against a
    repository}\label{feat:standards-engine}
  \featureitem{Deterministic audit}{per-standard checks (README badges, Python
    config, CI templates, shell scripts, naming conventions) that exit 0 or 1
    with structured findings}\label{feat:deterministic-audit}
  \featureitem{Preflight}{run quality gates and PII scan on staged files before
    commit --- the daily-use entry point}\label{feat:preflight}
  \featureitem{Custom standards directory}{teams point the engine at their own
    rules, not just punt-labs defaults}\label{feat:custom-standards}
  \featureitem{Claude Code plugin}{slash commands (\texttt{/punt audit},
    \texttt{/punt preflight}) that inject standards context into AI agent
    sessions}\label{feat:plugin}
  \featureitem{MCP server}{exposes audit, init, and preflight as MCP tools so
    any AI agent (Cursor, Windsurf, Codex, custom agents) can enforce the same
    standards without a plugin}\label{feat:mcp}
\end{enumerate}

\subsection*{Should Do}

Features that meaningfully improve the product but are not launch-blocking.
Candidates for fast follow-up.

\begin{enumerate}[nosep,leftmargin=2.5em]
  \featureitem{Release automation}{bump version, tag, push, and monitor the
    full pipeline (build, TestPyPI, test-install, PyPI) in one
    command}\label{feat:release}
  \featureitem{Cross-repo rollout}{detect which repos need a standard change,
    preview the delta, apply, commit, push, and report across N
    repos}\label{feat:rollout}
  \featureitem{Project doctor}{health check that verifies MCP wiring, plugin
    installation, quality gate configuration, and beads
    initialization}\label{feat:doctor}
  \featureitem{Status dashboard}{combined view of beads in-progress, git state,
    CI status, and unreleased commits}\label{feat:status}
  \featureitem{REST API}{HTTP interface for CI pipelines and internal dashboards
    to query audit findings programmatically}\label{feat:rest-api}
\end{enumerate}

\subsection*{Won't Do}

Features explicitly excluded. Naming what we will not build prevents scope
creep and clarifies the product's identity.

\begin{enumerate}[nosep,leftmargin=2.5em]
  \featureitem{GUI or web dashboard}{punt-kit is a terminal tool and a Claude
    Code plugin. A web interface would split focus and target a different user.
    Teams that want a portal should use
    Backstage}\label{feat:no-gui}
  \featureitem{CI system replacement}{punt-kit generates and validates CI
    workflows. It does not replace GitHub Actions, GitLab CI, or any CI
    runner}\label{feat:no-ci}
  \featureitem{Code parsing or linting}{punt-kit orchestrates linters (ruff,
    mypy, shellcheck) and validates their configuration. It does not parse
    source code. Teams that want a metalinter should use Trunk or
    MegaLinter}\label{feat:no-linting}
  \featureitem{Package management}{validates distribution config and automates
    release pipelines; does not replace PyPI, npm, or
    uv}\label{feat:no-package-mgmt}
\end{enumerate}

% ══════════════════════════════════════════════════════════════════════════════
% REFERENCES
% ══════════════════════════════════════════════════════════════════════════════

\newpage
\printbibliography[title={References}]

\end{document}
